% !PRINTED 88
\weberpage{88}
%<<CONT
When the imaginary part of $\omega$ grows to infinity, so
that the absolute value of $q$ vanishes, then
\[
\vartheta_{00}(v) = 1, \quad \vartheta_{01}(v) = 1, \quad q^{-\frac{1}{4}}\vartheta_{10}(v) = 2\cos\pi v,
\]
\[
q^{-\frac{1}{4}}\vartheta_{11}(v) = 2\sin\pi v.
\]
If we take $\omega$ purely imaginary, and hence $q$ real, positive and properly
fractional, then for a real $v$:

1. $\vartheta_{00}(v)$, $\vartheta_{01}(v)$, $\vartheta_{10}(v)$, $\vartheta_{11}(v)$ real, and when $v$ lies between 0
and $\frac{1}{2}$, positive [by § 24 (7)], consequently also $\vartheta_{00}$, $\vartheta_{01}$, $\vartheta_{10}$, $\vartheta'_{11}$
positive;

2. $\vartheta_{00}(iv)$, $\vartheta_{01}(iv)$, $\vartheta_{10}(iv)$, $-i\vartheta_{11}(iv)$ real, and as long as $v$
lies between 0 and $\dfrac{-i\omega}{2}$, positive; further, with the aid of the
formulae § 21, (8)

3. $\vartheta_{00}(\frac{1}{2}+iv)$, $\vartheta_{01}(\frac{1}{2}+iv)$, $i\vartheta_{10}(\frac{1}{2}+iv)$, $\vartheta_{11}(\frac{1}{2}+iv)$
real, and as long as $v$ lies between 0 and $\dfrac{-i\omega}{2}$, positive;
\[
4. \quad q^{\frac{1}{4}}e^{\pi iv}\vartheta_{00}\left(\frac{\omega}{2}+v\right), \quad -iq^{\frac{1}{4}}e^{\pi iv}\vartheta_{01}\left(\frac{\omega}{2}+v\right),
\]
\[
q^{\frac{1}{4}}e^{\pi iv}\vartheta_{10}\left(\frac{\omega}{2}+v\right), \quad -iq^{\frac{1}{4}}e^{\pi iv}\vartheta_{11}\left(\frac{\omega}{2}+v\right)
\]
real, and as long as $v$ lies between 0 and $\frac{1}{2}$, positive.

\begin{center}
\textbf{§ 26. \gesperrt{Development of $\vartheta$-Quotients.}}
\end{center}

If $v$ denotes a variable and $a$ an arbitrary constant, then
the quotients
\[
\frac{\vartheta_{g_1,g_2}(v+a)}{\vartheta_{g'_1,g'_2}(v)}
\]
are doubly periodic functions of the second kind and of the first order,
with the periodicity factors $(-1)^{g_1+g'_1}$, $(-1)^{g_2+g'_2}e^{-2\pi ia}$; and
if one adds an exponential factor $e^{\lambda v}$, then $\lambda$
and $a$ can be so determined that the periodicity factors become arbitrarily given quantities. By choosing the principal characteristics
$(g_1,g_2), (g'_1,g'_2)$ in all possible ways, one obtains 16 such
functions. We start out from one among them, for which, with the
addition of a constant factor, we
\weq{1}{F = \frac{\vartheta'_{11}\vartheta_{11}(v+a)}{2i\pi\vartheta_{11}(v)\vartheta_{11}(a)}}
%CONT>>
